%=============================================================================!
% 9.S  WHAT WE NOW HAVE                                                       !
%=============================================================================!
\unnumbered{What we now have}

A method of computing motion is a step map, and its order says how fast
its error over a fixed time shrinks with the step: the forward Euler
method, whose error over a fixed time shrinks as $\dt$ and over one step
as $\dt^2$, is of order 1. Adding the Taylor
series a step ahead and a step behind gives the Verlet method, which in
its velocity form is a half kick, a drift and a half kick, costs one force
evaluation per step and is of order 2; its local error, of order
$\dt^3$, is the `cubic scaling' quoted by TrajCast.

Reversing the
velocities after a step and stepping again returns the start, with its
velocities reversed, without error, and each step keeps
volume in phase space and satisfies the full symplectic condition. The
same method follows from splitting the Liouville operator into the parts
that move positions and momenta, whose exponentials are a drift and a kick
with no error; the symmetric product agrees with the true motion to second
order.

For a spring, velocity Verlet keeps the shadow energy $p^2/(2m) +
\frac12 kq^2(1 - \omega^2\dt^2/4)$ exactly, so the true energy oscillates
in a band of relative width $\omega^2\dt^2/4$; the symplectic Euler
method keeps one that differs from the energy at first order in $\dt$,
hence its wider band, and for other Hamiltonians a nearby shadow energy
is kept very nearly for sufficiently small steps and smooth forces.

The method is stable only for
$\omega\dt < 2$ for the fastest vibration, \SI{2.83}{\femto\second} for
the O--H stretch, and accurate enough well inside that; the step is chosen
by measuring the energy fluctuation. Over long runs it keeps the energy
better than RK4 at the same cost, provided the forces are smooth.

Chapter~10 uses velocity Verlet to build a complete simulation: atoms in a
periodic box, neighbour lists for the cutoffs, and the Ewald sum for
charges.
